\documentclass[a4paper,12pt]{letter}

% --- Set up fonts to use ---
\usepackage[bitstream-charter]{mathdesign}
\DeclareSymbolFont{usualmathcal}{OMS}{cmsy}{m}{n}
\DeclareSymbolFontAlphabet{\mathcal}{usualmathcal}
\usepackage{fontspec}
\setmainfont[Scale=0.9, Ligatures=TeX]{Charis SIL}
\defaultfontfeatures[Charis SIL]{Script=latn}
%\usepackage{microtype}                                    % see README
% --- END code for fonts ---
% --- Additional packages here ---
% Languages et al.
\usepackage{polyglossia}
\setmainlanguage{english}
\setkeys{english}{variant=british}
% \setmainlanguage{portuges}
% END Languages et al.

\usepackage[bottom]{footmisc}                               % make sure footnotes are placed on the bottom
\usepackage{endnotes}                                       % comment this ***and the next line*** if not needed
\renewcommand{\theendnote}{\Roman{endnote}}                 % Endnotes in Roman numerals (arabic if for footnotes)
\usepackage[absolute]{textpos}                              % For the positioning of the logo
% --- END additional packages ---
% --- Page layout ---
\usepackage[top=2.5cm, bottom=2.5cm, left=3.0cm, right=3.0cm]{geometry}
\pdfpagewidth 210mm
\pdfpageheight 297mm
% no line numbering
\pagestyle{empty}
% avoid having the closing too ``centered''
\longindentation=230pt
% --- END Page layout ---
% --- Commands (new or redefine) ---
% --- END commands ---
% --- Includes ---
% general purpose
\newcommand{\ts}{\textsection}
\newcommand{\emd}{\textemdash}                             % see [1]
\newcommand{\fn}[1]{\footnote{#1}}
\renewcommand{\tt}[1]{\texttt{#1}}

%% prevent figures that end up on the last empty 
%% page of a document from being forcibly centred 
%% on the middle of the page.
\makeatletter
\setlength{\@fptop}{0pt}
\makeatother

% [1] - When inserting an em-dash like this: --- any adjacent words 
%       are NOT HYPHENATED. Using \textemdash (or the newly defined 
%       shortcut)fixes this issue.

%% --- math packages (in ArchLinux they're all part of TeXLive) ---

\usepackage{amsmath}   % has to come before eulervm, or things bork!
\usepackage{amsthm}    % required for proof environ (et al.)

% --- END math packages ---

% math stuff
\newcommand{\eqv}{\Leftrightarrow}
\newcommand{\pvec}[1]{\overrightarrow{#1}}
\newcommand{\npvec}[1]{\norm{\pvec{#1}}}
\newcommand{\mb}[1]{\ensuremath \mathbb{#1}}
\newcommand{\mc}[1]{\ensuremath \mathcal{#1}}
\newcommand{\f}[2]{\ensuremath \frac{#1}{#2}}

\providecommand{\abs}[1]{\left \lvert #1 \right \rvert}
\providecommand{\norm}[1]{\left \lVert #1 \right \rVert}

\providecommand{\pth}[1]{\left ( #1 \right ) }
\providecommand{\spth}[1]{\left [ #1 \right ] }
\providecommand{\cpth}[1]{\left { #1 \right } }

% make the optional theorem name bold
\makeatletter
\def\th@plain{%
  \thm@notefont{}% same as heading font
  \itshape % body font
}
\makeatother

% math theorem environments
\newtheorem{thm}{Theorem}[section]
\newtheorem{corollary}[thm]{Corollary}
\newtheorem{lemma}[thm]{Lemma}
\newtheorem{definition}[thm]{Definition} % 'def' cannot be used as environ name

\theoremstyle{remark}
\newtheorem{remark}[thm]{Remark}

\renewcommand\qedsymbol{$\blacksquare$}  % for proofs to end with a black *filled* square (requires amssymb)

%
% math hacks
%

% scalable middle bar; see 
% http://tex.stackexchange.com/questions/5502/how-to-get-a-mid-binary-relation-that-grows
\newcommand{\relmiddle}[1]{\mathrel{}\middle#1\mathrel{}}

% it's a cheat but this makes spacing less wrong around quantifiers
\DeclareMathOperator{\Nexists}{\nexists}
\DeclareMathOperator{\Exists}{\exists}
\DeclareMathOperator{\Forall}{\forall}

% --- END Includes ---

% --- Letter specific info ---
% NOTA BENE:
% 1) Stripping \address makes closing and signature *left aligned*
% 2) Stripping \opening removes both opening AND date & address from top of letter

% for letters, I prefer to set the date explicitly
\date{February 31, 2013} 
\address{
  22nd Baking Street\\
  2099--101 Acme City\\
  LooneyTuneLand
}
\signature{Your Truly's Name}
% --- END PREAMBLE ---

\setlength{\TPHorizModule}{30mm}
\setlength{\TPVertModule}{\TPHorizModule}

\begin{document}
\begin{letter}{
	R. MaDillo\\
	European Parliament\\
	Rue Wiertz\\
	Altiero Spinelli 07F343\\
	B-1047 Bruxelles 
}

\begin{textblock}{10}(1.0, 0.3)
	\includegraphics{logo}
\end{textblock}

%\opening{Ex.\textsuperscript{o} Doutor R. Madillo,} % -- Portuguese introduction style
\opening{Dear Dr. R. MaDillo,}
{
\begin{center}
	\textbf{\textsc{Subject of your letter here}}\\[2em]
\end{center}
}

% 
% By default, the paragraphs are separated by a blank line, and there is no 
% first line indentation. When writing in Portuguese, it is customary to do the 
% *opposite*, i.e. indent the paragraph's first line and don't use blank lines 
% between paragraphs. To obtain this setup, UNcomment the two \setlength and the 
% \vspace commands below, as well as the two \setlength before \closing.
%
% \setlength{\parindent}{0.5cm}
% \setlength{\parskip}{0pt plus 0.1pt}
% \vspace{1em}
%

	Lorem ipsum dolor sit amet, consectetur adipiscing elit. Donec bibendum 
	suscipit odio, id dapibus ipsum placerat vitae. Quisque vitae magna vel ante 
	semper dapibus sit amet eu elit. Fusce cursus sodales eros, non porta urna 
	bibendum et. Ut mauris risus, malesuada sit amet semper ac, molestie in 
	tellus. Nam imperdiet metus sed massa faucibus sed vulputate risus laoreet. 
	Aenean tincidunt, purus ac euismod lobortis, tortor lacus iaculis eros, non 
	egestas orci massa in ligula. Cras bibendum suscipit lacus nec mollis. Nullam 
	neque erat, molestie ac bibendum eget, molestie non nunc. Morbi quis orci 
	augue. Pellentesque facilisis, erat sed pellentesque rutrum, orci turpis 
	dictum mi, sit amet rutrum nibh quam in velit. Vestibulum semper, est eget 
	blandit pellentesque, nisi elit venenatis diam, id vestibulum velit leo quis 
	orci.\endnote{Some endnote-ing text.} 

	In dui magna, posuere eget, vestibulum et, tempor auctor, justo. In ac felis 
	quis tortor malesuada pretium. Pellentesque auctor neque nec urna. Proin 
	sapien ipsum, porta a, auctor quis, euismod ut, mi. Aenean viverra rhoncus 
	pede. 

% \setlength{\parindent}{0.0cm}
% \setlength{\parskip}{0.7em}
\closing{Yours sincerely,}
%
\ps{P.\ S.: Some more latin should go here.}
\encl{A vacation postcard!}

% --- show endnotes ---
% comment if not using the endnotes package
\vspace{1cm}
\begingroup
%\addcontentsline{toc}{section}{Notes}
\def\enotesize{\small}
\def\enoteheading{\Large \textsc{Notes}}
\theendnotes
\endgroup
% --- end endnotes ---
\end{letter}
\end{document}

% ALTERNATIVE CLOSING: if you want to send a letter that will not be 
% printed, you can add an image with your signature replacing the \closing 
% with this:
%
%  \closing{Yours sincerely,\\
%    \fromsig{\includegraphics[scale=0.5]{signature.png}}\\
%    \fromname{Your Truly's Name}
%  }
%
% NOTE: If you use this, you have to remove the \signature, in the preamble.
